% GENERATED FILE. Edit canonical lessons, then run npm run generate:books.
% canonical-source-hash: fnv1a64:c48851712f41b4e8
% canonical-lessons: ML-C345-duhkham, ML-W345-visarga, ML-C345-santosham-tonni, ML-C345-kanan-agrahamundu, ML-R345-feelings-and-hopes

\chapter{Feelings and Hopes}
\label{ch:ml-feelings-and-hopes}
\hlchaptermodality{345}

\begin{chapteropening}
\textbf{By the end of this chapter:} \emph{I can say how I felt and how I feel with the dative and tōnnuka, tell a feeling as scene or as event, and say what I hope to do or become.}

Complete ML-R345-feelings-and-hopes: feelings with enikkŭ … tōnni, sorrow with duḥkhaṁ and its visarga, and a wish with -ān āgrahamuṇṭŭ.
\end{chapteropening}

\section[\texorpdfstring{\ml{ദുഃഖം} (duḥkhaṁ)}{ദുഃഖം (duḥkhaṁ)}]{\ml{ദുഃഖം} (duḥkhaṁ) — grief, sorrow}

\label{lesson:ML-C345-duhkham}



\begin{quote}
Before the new one: say \textquotedblleft{}sadness\textquotedblright{}, then \textquotedblleft{}well-being\textquotedblright{}, then \textquotedblleft{}news\textquotedblright{}. (\emph{saṅkaṭaṁ}, \emph{sukhaṁ}, \emph{vārtta}.)
\end{quote}

\subsection*{You'll want to know: \ml{ദുഃഖം}}
\textbf{\ml{ദുഃഖം}} — \emph{duḥkhaṁ} — \textquotedblleft{}grief, sorrow\textquotedblright{}.

\begin{quote}
\textbf{\ml{ആ} \ml{വാർത്ത} \ml{കേട്ടപ്പോൾ} \ml{എനിക്ക്} \ml{വലിയ} \ml{ദുഃഖമുണ്ടായി}.}
\end{quote}

\begin{quote}
\emph{ā vārtta kēṭṭappōḷ enikkŭ valiya duḥkhamuṇṭāyi.} — \textquotedblleft{}When I heard that news, I felt deep sorrow.\textquotedblright{}
\end{quote}

\textbf{\ml{സങ്കടം}} is the everyday sadness of a bad day. \textbf{\ml{ദുഃഖം}} is heavier: grief, or sorrow you speak of with care.

\begin{quote}
\textbf{\ml{അവരുടെ} \ml{ദുഃഖം} \ml{എനിക്ക്} \ml{മനസ്സിലാകുന്നു}.}
\end{quote}

\begin{quote}
\emph{avaruṭe duḥkhaṁ enikkŭ manassilākunnu.} — \textquotedblleft{}I understand their grief.\textquotedblright{}
\end{quote}

The mark after \textbf{\ml{ദു}} is a breath, \emph{ḥ}. Many speakers let it turn into a held \emph{k}, close to \emph{dukkhaṁ}. The next lesson puts the mark in your hand.

\begin{cousinweb}
\textbf{\ml{ദുഃഖം}} (\emph{duḥkhaṁ}) $\leftarrow$ Sanskrit \emph{duḥkha}, \textquotedblleft{}difficulty, pain\textquotedblright{}, a tatsama with Malayalam's noun ending (CDIAL 6375). CDIAL notes it was formed after \emph{sukha}, the word you know as \textbf{\ml{സുഖം}}: the two are a pair, ease and its opposite.
\end{cousinweb}

\subsection*{Your turn}
Say these aloud:

\begin{itemize}
\raggedright
  \item \emph{duḥkhaṁ}, with a short breath after \emph{du}
  \item \textquotedblleft{}When I heard that news, I felt deep sorrow.\textquotedblright{}
  \item the everyday word for sadness, then the heavier one
  \item the pair: \emph{sukhaṁ}, \emph{duḥkhaṁ}
\end{itemize}

\begin{tcolorbox}[breakable,colback=teal!4,colframe=teal!35!black,title={Before you move on}]
What does \textbf{\ml{ദുഃഖം}} mean? (\textbf{Grief, sorrow.}) How is it different from \textbf{\ml{സങ്കടം}}? (\textbf{it is heavier.}) Which word is its opposite? (\textbf{\ml{സുഖം}}.)
\end{tcolorbox}
\section[\texorpdfstring{\ml{ഃ} (ḥ)}{ഃ (ḥ)}]{\ml{ഃ} (ḥ) — the two rings inside \ml{ദുഃഖം}}

\label{lesson:ML-W345-visarga}



\begin{quote}
Before the new one: write \textbf{\ml{ം}} and \textbf{\ml{ഖ}} once each, and name them. (The anusvaram, \emph{ṁ}, and \emph{kha}.)

Now say \textbf{\ml{ദുഃഖം}} (\emph{duḥkhaṁ}), \textquotedblleft{}grief\textquotedblright{}. Both of those are in it. This lesson takes out the one mark you have not written and puts it in your hand.
\end{quote}

\subsection*{Script you'll notice: \ml{ഃ}}
\textbf{\ml{ഃ}} — the visarga, \emph{ḥ}: a short breath after a vowel. It is two small rings, one above the other, written after its syllable: \textbf{\ml{ദു}} + \textbf{\ml{ഃ}} $\to$ \textbf{\ml{ദുഃ}}, \emph{duḥ}.

Do not confuse it with \textbf{\ml{ം}}, one ring, which is \emph{ṁ}. The visarga lives in words taken from Sanskrit, and most of them are learned words.

\subsection*{Writing: \ml{ഃ}}
One recorded hand writes \textbf{\ml{ഃ}} in two strokes: the upper ring first, then, after a lift, the lower one. Each ring starts at its top and goes round anticlockwise, to the left first.

Write \textbf{\ml{ഃ}} yourself, slowly, and larger than it is printed: two rings the same size, one above the other. Copy the printed shape and keep the order above.

\subsection*{Your turn}
\begin{itemize}
\raggedright
  \item \emph{Find:} \ml{ഃ} inside \textbf{\ml{ദുഃഖം}}
  \item \emph{Trace it:} \ml{ഃ} three times, saying \emph{ḥ} as you finish each one
  \item \emph{Write it:} \ml{ം}, then \ml{ഃ}, from memory, naming each as you finish it
  \item \emph{Write it:} \textbf{\ml{ദുഃഖം}}
\end{itemize}

\begin{tcolorbox}[breakable,colback=teal!4,colframe=teal!35!black,title={Before you move on}]
Which mark is this — \ml{ഃ}? (\textbf{the visarga}, \emph{ḥ}, a breath.) Which word did it come from? (\textbf{\ml{ദുഃഖം}}, \emph{duḥkhaṁ}, \textquotedblleft{}grief\textquotedblright{}.) How is it different from \textbf{\ml{ം}}? (\textbf{Two rings, not one.})
\end{tcolorbox}
\section[\texorpdfstring{\ml{എനിക്ക്} \ml{സന്തോഷം} \ml{തോന്നി} (enikkŭ santōṣaṁ …)}{എനിക്ക് സന്തോഷം തോന്നി (enikkŭ santōṣaṁ …)}]{\ml{എനിക്ക്} \ml{സന്തോഷം} \ml{തോന്നി} (enikkŭ santōṣaṁ tōnni) — I felt happy}

\label{lesson:ML-C345-santosham-tonni}



\begin{quote}
Say \textquotedblleft{}joy\textquotedblright{}, \textquotedblleft{}sadness\textquotedblright{}, \textquotedblleft{}fear\textquotedblright{}, and \textquotedblleft{}tiredness\textquotedblright{}. (\emph{santōṣaṁ}, \emph{saṅkaṭaṁ}, \emph{pēṭi}, \emph{kṣīṇaṁ}.)
\end{quote}

\begin{grammarlens}[title={Grammar Lens: the feeling comes to you}]
English says \textquotedblleft{}I felt happy\textquotedblright{}, with \textquotedblleft{}I\textquotedblright{} in charge. Malayalam puts you in the dative and lets the feeling arrive: \textquotedblleft{}to me, joy seemed\textquotedblright{}.

\begin{quote}
\textbf{\ml{വാർത്ത} \ml{കേട്ടപ്പോൾ} \ml{എനിക്ക്} \ml{സന്തോഷം} \ml{തോന്നി}.}
\end{quote}

\begin{quote}
\emph{vārtta kēṭṭappōḷ enikkŭ santōṣaṁ tōnni.} — \textquotedblleft{}When I heard the news, I felt happy.\textquotedblright{}
\end{quote}

\begin{quote}
\textbf{\ml{എനിക്ക്} \ml{ക്ഷീണം} \ml{തോന്നുന്നു}.}
\end{quote}

\begin{quote}
\emph{enikkŭ kṣīṇaṁ tōnnunnu.} — \textquotedblleft{}I feel tired.\textquotedblright{}
\end{quote}

The noun names the feeling; \textbf{\ml{തോന്നി}}, \textbf{\ml{തോന്നുന്നു}} carry the time. \textbf{\ml{അവന്}}, \textbf{\ml{അവൾക്ക്}}, \textbf{\ml{ഞങ്ങൾക്ക്}} take the same place for \textquotedblleft{}he\textquotedblright{}, \textquotedblleft{}she\textquotedblright{}, \textquotedblleft{}we\textquotedblright{}.
\end{grammarlens}

\begin{grammarlens}[title={Grammar Lens: a feeling as scene, or as event}]
A feeling that lasted is scene, in \textbf{\ml{ആയിരുന്നു}}. A feeling that struck is an event, and Malayalam says it came:

\begin{quote}
\textbf{\ml{ഇന്നലെ} \ml{എനിക്ക്} \ml{വലിയ} \ml{സങ്കടമായിരുന്നു}.}
\end{quote}

\begin{quote}
\emph{innale enikkŭ valiya saṅkaṭamāyirunnu.} — \textquotedblleft{}Yesterday I was very sad.\textquotedblright{}
\end{quote}

\begin{quote}
\textbf{\ml{അത്} \ml{കേട്ടപ്പോൾ} \ml{അവന്} \ml{ദേഷ്യം} \ml{വന്നു}.}
\end{quote}

\begin{quote}
\emph{atŭ kēṭṭappōḷ avanŭ dēṣyaṁ vannu.} — \textquotedblleft{}When he heard that, he got angry.\textquotedblright{}
\end{quote}

To describe a person, add \textbf{-\ml{ഉള്ള}}: \textbf{\ml{സന്തോഷമുള്ള} \ml{കുട്ടി}} (\emph{santōṣamuḷḷa kuṭṭi}), \textquotedblleft{}a happy child\textquotedblright{}.
\end{grammarlens}

\subsection*{Your turn}
Say these aloud:

\begin{itemize}
\raggedright
  \item \textquotedblleft{}When I heard the news, I felt happy.\textquotedblright{}
  \item \textquotedblleft{}I feel tired\textquotedblright{}, then \textquotedblleft{}I felt afraid\textquotedblright{} — \emph{enikkŭ pēṭi tōnni}
  \item \textquotedblleft{}Yesterday I was very sad\textquotedblright{}, then \textquotedblleft{}he got angry\textquotedblright{}
  \item \textquotedblleft{}a happy child\textquotedblright{}
\end{itemize}

\begin{tcolorbox}[breakable,colback=teal!4,colframe=teal!35!black,title={Before you move on}]
\textquotedblleft{}I felt happy\textquotedblright{}? (\emph{\textbf{enikkŭ santōṣaṁ tōnni.}}) A feeling that lasted? (\emph{\textbf{āyirunnu}}: \emph{enikkŭ saṅkaṭamāyirunnu}.) A feeling that struck? (It \textbf{came}: \emph{avanŭ dēṣyaṁ vannu}.)
\end{tcolorbox}
\section[\texorpdfstring{\ml{എനിക്ക്} \ml{കേരളം} … (enikkŭ kēraḷaṁ …)}{എനിക്ക് കേരളം … (enikkŭ kēraḷaṁ …)}]{\ml{എനിക്ക്} \ml{കേരളം} \ml{കാണാൻ} \ml{ആഗ്രഹമുണ്ട്} (enikkŭ kēraḷaṁ kāṇān āgrahamuṇṭŭ) — I would like to see Kerala}

\label{lesson:ML-C345-kanan-agrahamundu}



\begin{quote}
Say \textquotedblleft{}a wish\textquotedblright{}, \textquotedblleft{}a dream\textquotedblright{}, and \textquotedblleft{}to read\textquotedblright{} in the \textbf{-\ml{ആൻ}} form. (\emph{āgrahaṁ}, \emph{svapnaṁ}, \emph{vāyikkān}.)
\end{quote}

\begin{grammarlens}[title={Grammar Lens: -\ml{ആൻ} \ml{ആഗ്രഹമുണ്ട്}}]
The \textbf{-\ml{ആൻ}} form you already know says \textquotedblleft{}to do it\textquotedblright{}. Put it before \textbf{\ml{ആഗ്രഹം}}, \textquotedblleft{}a wish\textquotedblright{}, and the wish belongs to someone in the dative, with \textbf{\ml{ഉണ്ട്}}:

\begin{quote}
\textbf{\ml{എനിക്ക്} \ml{കേരളം} \ml{കാണാൻ} \ml{ആഗ്രഹമുണ്ട്}.}
\end{quote}

\begin{quote}
\emph{enikkŭ kēraḷaṁ kāṇān āgrahamuṇṭŭ.} — \textquotedblleft{}I would like to see Kerala.\textquotedblright{}
\end{quote}

The verb \textbf{\ml{ആഗ്രഹിക്കുക}}, \textquotedblleft{}to wish, to long\textquotedblright{}, says the same with \textbf{\ml{ഞാൻ}} in front:

\begin{quote}
\textbf{\ml{ഞാൻ} \ml{ഒരു} \ml{അധ്യാപകൻ} \ml{ആകാൻ} \ml{ആഗ്രഹിക്കുന്നു}.}
\end{quote}

\begin{quote}
\emph{ñān oru adhyāpakan ākān āgrahikkunnu.} — \textquotedblleft{}I want to become a teacher.\textquotedblright{}
\end{quote}

\textbf{\ml{ആകുക}} is \textquotedblleft{}to become\textquotedblright{}; \textbf{\ml{ആകാൻ}}, \textquotedblleft{}to become\textquotedblright{}.
\end{grammarlens}

\begin{grammarlens}[title={Grammar Lens: … \ml{എന്നാണ്} \ml{എന്റെ} \ml{ആഗ്രഹം}}]
For a bigger hope, say the hope first, close it with \textbf{\ml{എന്നാണ്}}, \textquotedblleft{}is that …\textquotedblright{}, and name it last:

\begin{quote}
\textbf{\ml{ഒരു} \ml{ചെറിയ} \ml{കട} \ml{തുടങ്ങണം} \ml{എന്നാണ്} \ml{എന്റെ} \ml{ആഗ്രഹം}.}
\end{quote}

\begin{quote}
\emph{oru ceṟiya kaṭa tuṭaṅṅaṇaṁ ennāṇŭ enṟe āgrahaṁ.} — \textquotedblleft{}My wish is to start a small shop.\textquotedblright{}
\end{quote}

\textbf{\ml{ആഗ്രഹം}} can give way to \textbf{\ml{സ്വപ്നം}}, \textquotedblleft{}dream\textquotedblright{}, or \textbf{\ml{ലക്ഷ്യം}}, \textquotedblleft{}aim\textquotedblright{}, in the same frame.
\end{grammarlens}

\subsection*{Your turn}
Say these aloud:

\begin{itemize}
\raggedright
  \item \textquotedblleft{}I would like to see Kerala.\textquotedblright{}
  \item \textquotedblleft{}I want to become a teacher.\textquotedblright{}
  \item \textquotedblleft{}My wish is to start a small shop\textquotedblright{}, then the same with \emph{svapnaṁ}
  \item one thing you hope to do next year, and why
\end{itemize}

\begin{tcolorbox}[breakable,colback=teal!4,colframe=teal!35!black,title={Before you move on}]
\textquotedblleft{}I would like to see Kerala\textquotedblright{}? (\emph{\textbf{enikkŭ kēraḷaṁ kāṇān āgrahamuṇṭŭ.}}) \textquotedblleft{}To become\textquotedblright{}? (\emph{\textbf{ākān.}}) \textquotedblleft{}My aim is to …\textquotedblright{}? (\textbf{… -\ml{അണം} \ml{എന്നാണ്} \ml{എന്റെ} \ml{ലക്ഷ്യം}}.)
\end{tcolorbox}
\section[(review)]{Review — feelings and hopes}

\label{lesson:ML-R345-feelings-and-hopes}



\begin{quote}
Say \textquotedblleft{}grief, sorrow\textquotedblright{}. (\emph{duḥkhaṁ}.) Write \textbf{\ml{ഃ}} once from memory, then \textbf{\ml{ദുഃഖം}}.
\end{quote}

\subsection*{Your turn}
Say these aloud:

\begin{itemize}
\raggedright
  \item \textquotedblleft{}When I heard the news, I felt happy.\textquotedblright{}
  \item \textquotedblleft{}When I heard that news, I felt deep sorrow.\textquotedblright{}
  \item \textquotedblleft{}Yesterday I was very sad. Then he got angry.\textquotedblright{}
  \item \textquotedblleft{}I would like to see Kerala. I want to become a teacher.\textquotedblright{}
  \item \textquotedblleft{}My wish is to start a small shop.\textquotedblright{}
  \item \textquotedblleft{}My bag got lost, and I felt very sad.\textquotedblright{} — \emph{enṟe bāgŭ naṣṭappeṭṭupōyi, enikkŭ valiya saṅkaṭaṁ tōnni.}
  \item \textquotedblleft{}Is tomorrow convenient for you?\textquotedblright{} — \emph{nāḷe niṅṅaḷkkŭ saukaryamāṇō?}
\end{itemize}

\begin{tcolorbox}[breakable,colback=teal!4,colframe=teal!35!black,title={Before you move on}]
\textquotedblleft{}I felt happy\textquotedblright{}? (\emph{\textbf{enikkŭ santōṣaṁ tōnni.}}) \textquotedblleft{}He got angry\textquotedblright{}? (\emph{\textbf{avanŭ dēṣyaṁ vannu.}}) \textquotedblleft{}I would like to see Kerala\textquotedblright{}? (\emph{\textbf{enikkŭ kēraḷaṁ kāṇān āgrahamuṇṭŭ.}}) \textquotedblleft{}Grief\textquotedblright{}? (\textbf{\ml{ദുഃഖം}}, with the visarga \textbf{\ml{ഃ}}.)
\end{tcolorbox}
